\chapter{Encours vendeur à découvert, emprunt et encombrement}\label{ch:s1:short-interest-borrow-and-crowding}

En janvier 2021, l'encours vendeur à découvert de GameStop atteignait 122.97~\% de son flottant~: plus d'actions avaient été vendues à découvert qu'il n'en était disponible à la négociation, parce que les
mêmes actions avaient été prêtées, vendues puis de nouveau prêtées. Le 27 janvier, le titre clôturait à 347.51~\$, soit plus de 1\,600~\% au-dessus de son cours de clôture seize jours plus tôt. Le personnel de la SEC
a constaté ensuite que les vendeurs à découvert qui se couvraient avaient contribué à certaines des hausses les plus fortes, bien que leurs achats ne représentent qu'une petite fraction de
l'ensemble des achats. L'encours vendeur à découvert est généralement lu comme un signal, car les vendeurs à découvert ont tendance à être informés~; dans un titre encombré, c'est aussi une position qui peut être
dénouée de force. Ce chapitre construit un marché de prêt de titres sur les titres synthétiques, mesure ce que prédisent l'encours vendeur à découvert et les frais d'emprunt, ce que coûtent ces frais, et les pertes possibles
d'un portefeuille vendeur encombré. La construction est \texttt{firm.lendsig}.

\section{L'encours vendeur à découvert comme signal}

L'encours vendeur à découvert (livre~1, chapitre~16) est le nombre d'actions vendues à découvert et non encore rachetées, exprimé en part des actions en circulation ou du flottant,
ou en jours de volume moyen (jours de couverture). Les éléments montrant qu'il prédit les rendements sont solides sous certaines formes et fragiles sous d'autres. Boehmer, Jones et Zhang,
à partir de données quotidiennes de vente à découvert du NYSE pour 2000--2004, ont constaté que les titres fortement vendus à découvert sous-performaient les titres peu vendus d'un écart ajusté du risque de 1.16~\% sur les
vingt jours de bourse suivants~: les vendeurs à découvert, pris collectivement, sont informés. Asquith, Pathak et Ritter, à partir de l'encours mensuel sur une période plus longue, ont constaté que bon nombre des
régularités documentées n'étaient pas robustes~: les portefeuilles d'actions à fort encours vendeur, pondérés également, ont généralement sous-performé, alors que ceux pondérés par la capitalisation ne l'ont pas été.

\begin{dated}{2026-09}{Déclaration de l'encours vendeur aux États-Unis}\label{dat:s1:short-interest-borrow-and-crowding:finra}
En vertu de la FINRA Rule 4560, les broker-dealers déclarent les positions vendeuses brutes sur l'ensemble des titres de capital détenus dans leurs comptes clients et pour compte propre deux fois par
mois, aux dates de règlement-livraison désignées par la FINRA~; les déclarations doivent parvenir avant 18~h, heure de l'Est, le deuxième jour ouvré suivant la date de règlement-livraison, et la FINRA
publie l'encours vendeur agrégé selon un calendrier publié, quelques jours plus tard. Un signal fondé sur l'encours voit donc les positions telles qu'elles étaient une à
plusieurs semaines plus tôt~; les volumes quotidiens de vente à découvert et les données de prêt de titres sont les alternatives commerciales, plus rapides.
\end{dated}

Le marché synthétique ne compte aucun vendeur à découvert~: \texttt{firm.lendsig} les ajoute donc. Chaque titre dispose d'une offre prêtable, égale à une part lognormale de ses actions, centrée sur
20~\%. La demande vendeuse a deux composantes~: les vendeurs informés, qui suivent le rendement espéré injecté dans le titre avec un décalage de vingt jours et vendent en proportion de
son caractère négatif~; et les vendeurs non informés, bruit persistant d'une médiane de 2~\% des actions. L'encours vendeur est la demande que l'offre peut satisfaire. Sur les années 3 à
10, la médiane des titres a 3.1~\% de ses actions vendues à découvert, le 90e centile 7.5~\% et le titre le plus vendu 54~\%~; la médiane des jours de couverture est de 2.1, et le 90e
centile 6.8. Parce que les vendeurs informés figurent par construction dans la demande, l'encours vendeur prédit les rendements~: l'opposé de l'encours vendeur présente un IC de rang mensuel de
0.061, les jours de couverture 0.048, et un score d'encombrement (le rang moyen de l'encours vendeur, du taux d'utilisation et des jours de couverture) 0.058.

\section{Frais d'emprunt et taux d'utilisation}

\begin{definition}[Signal du coût d'emprunt des titres]\label{def:s1:short-interest-borrow-and-crowding:fee}
Un \emph{signal du coût d'emprunt des titres}\index{signal du coût d'emprunt des titres} classe les titres selon les frais d'emprunt requis pour une vente à découvert (ou selon le taux d'utilisation, soit la part de l'offre prêtable
effectivement en prêt)~; des frais élevés signalent des titres soumis à une forte demande vendeuse rapportée à l'offre, qui ont eu tendance à sous-performer.
\end{definition}

Les frais d'emprunt synthétiques s'élèvent à 0.25~\% par an (collatéral général) tant que le taux d'utilisation n'atteint pas 60~\%, puis grimpent fortement jusqu'à 40~\% à pleine utilisation. Seuls 2.9~\% des
jours-titres paient plus de 1~\%, mais le centile supérieur paie le maximum. Les frais représentent le coût du signal et, ici, ils déterminent le résultat~:

\begin{center}
\small
\begin{tabular}{@{}lccc@{}}
\toprule
portefeuilles par décile, reconstitués chaque mois & encours vendeur & encombrement & frais d'emprunt\\
\midrule
IC de rang sur le mois suivant & 0.061 & 0.058 & 0.006\\
ratio de Sharpe avant frais & 3.33 & 3.30 & 1.18\\
ratio de Sharpe après frais & 2.24 & 2.08 & $-0.66$\\
ratio de Sharpe après frais et coûts & 1.84 & 1.69 & $-0.82$\\
rendement brut annuel & 8.3~\% & 8.2~\% & 2.5~\%\\
frais d'emprunt~; coûts de transaction (par an) & 2.7~\%~; 0.95~\% & 3.0~\%~; 0.94~\% & 4.0~\%~; 0.33~\%\\
frais moyens côté vente & 5.4~\% & 6.1~\% & 7.9~\%\\
\bottomrule
\end{tabular}
\end{center}

Le portefeuille fondé sur l'encours vendeur conserve l'essentiel de son avantage après frais (8.3~\% brut, 4.6~\% net de frais et de coûts). Le portefeuille fondé sur les frais sert d'avertissement~: les titres dont les frais sont les
plus élevés sous-performent, mais de moins que ce que coûte leur vente à découvert, et le portefeuille perd de l'argent. Engelberg, Reed et Ringgenberg ont formulé la même remarque sur des données réelles
que le niveau des frais n'est pas le seul risque~: les titres dont les prêts peuvent devenir coûteux ou être rappelés comportent un risque lié à la vente à découvert, et ont des rendements plus faibles, une moindre
efficience des cours et moins de ventes à découvert.

\begin{omfigure}
\begin{tikzpicture}
\begin{axis}[width=10cm,height=5.2cm,xlabel={\small année},ylabel={\small rendement cumulé (\%)},tick label style={font=\footnotesize},
  xmin=2,xmax=10,grid=major,grid style={gray!25},legend pos=north west,legend style={font=\scriptsize},table/col sep=comma]
\addplot[black!50,thick,dashed] table[x=year,y=gross]{figdata/strategies-1/08-short-interest-borrow-and-crowding/book.csv};
\addplot[omDef!60,thick] table[x=year,y=after_fee]{figdata/strategies-1/08-short-interest-borrow-and-crowding/book.csv};
\addplot[omDef,very thick] table[x=year,y=net]{figdata/strategies-1/08-short-interest-borrow-and-crowding/book.csv};
\legend{avant frais d'emprunt,après frais d'emprunt,après frais et coûts de transaction}
\end{axis}
\end{tikzpicture}

\omcaption{Le portefeuille par décile fondé sur l'encours vendeur sur le marché synthétique de prêt de titres~: rendement cumulé avant frais d'emprunt, après ceux-ci, et après coûts de transaction également
(exposition brute égale à un). Données~: \texttt{s1\_lending.book}.}\label{fig:s1:short-interest-borrow-and-crowding:book}
\end{omfigure}

L'usage le plus sûr du signal est l'évitement. Un portefeuille long seul, pondéré à parts égales, sur les titres synthétiques rapporte 9.1~\% par an~; le même portefeuille privé de son décile le plus
vendu rapporte 10.0~\%, soit un écart de 0.94~\% par an avec une statistique $t$ de 7.2, et il n'acquitte aucun frais.

\section{Rappels, squeezes et risque du marché de prêt}

\begin{definition}[Risque de squeeze]\label{def:s1:short-interest-borrow-and-crowding:squeeze}
Le \emph{risque de squeeze}\index{risque de squeeze} est le risque que les vendeurs à découvert d'un titre soient contraints de racheter au même moment (parce que des prêts sont rappelés,
qu'un appel de marge est lancé ou que les pertes s'accroissent), ce qui fait monter le cours et force à acheter davantage.
\end{definition}

Une position vendeuse présente trois risques qu'une position longue ne présente pas~: les frais peuvent augmenter, le prêt peut être rappelé et la perte est illimitée. Sur le marché synthétique,
les prêts sont rappelés avec une probabilité quotidienne de 0.05~\% à faible taux d'utilisation, qui monte à 2~\% à pleine utilisation~: les titres visés par un signal d'encours vendeur sont
ceux qui ont le plus de chances de voir leurs prêts retirés. Le squeeze est le cas extrême. Le personnel de la SEC a décrit son mécanisme pour GameStop~: des appels de marge ou des pertes poussent
les vendeurs à découvert à racheter, leur achat fait monter le cours, et la hausse oblige davantage de vendeurs à acheter~; ils ont aussi constaté que le cours est resté élevé après que les rachats
se sont essoufflés, de sorte que le squeeze n'explique qu'une partie de l'épisode, et non pas tout.

\section{Ventes à découvert encombrées}

\begin{definition}[Mesure d'encombrement vendeur]\label{def:s1:short-interest-borrow-and-crowding:crowding}
Une \emph{mesure d'encombrement vendeur}\index{mesure d'encombrement vendeur} évalue le nombre de vendeurs à découvert présents dans un titre au regard des issues dont ils disposent pour en sortir,
en combinant l'encours vendeur à découvert, le taux d'utilisation de l'offre prêtable et les jours de couverture.
\end{definition}

Le test de stress de squeeze (\texttt{firm.lendsig.squeeze\_loss}) constitue un portefeuille vendeur sur les vingt titres les plus
encombrés, pour la moitié de son capital au total, à chaque fin de mois, et demande ce qu'il perdrait si la moitié de l'encours vendeur de chaque titre est rachetée sur cinq jours, le
cours suivant la loi d'impact en racine carrée (livre~7, chapitre~27) tant que dure l'achat. La perte médiane représente 4.0~\% du capital et la pire 8.3~\%~;
la pire perte réelle sur cinq jours de ce même portefeuille sur les huit années est de 3.0~\%. Si seul un quart est couvert, la médiane tombe à 2.8~\%~; si la moitié est couverte sur dix
jours, 5.7~\% et au pire 12.3~\%. La loi en racine carrée n'a pas été conçue pour des achats portant sur de nombreux jours de volume, ces chiffres sont donc des planchers et non des prévisions~: le cours de GameStop
a progressé d'environ 2\,700~\% entre son point bas et son point haut. À cette échelle, l'arithmétique est plus simple que n'importe quel modèle~: un titre qui se multiplie par vingt dans un portefeuille vendeur sur vingt titres
à parts égales, avec la moitié de son capital, coûte 47.5~\% du capital, à lui seul.

\begin{omfigure}
\begin{tikzpicture}
\begin{axis}[ybar,bar width=6pt,width=11cm,height=5.0cm,xlabel={\small rendement sur cinq jours du portefeuille vendeur encombré (\% du capital)},
  ylabel={\small fins de mois},tick label style={font=\footnotesize},xmin=-10,xmax=3.5,ymin=0,grid=major,grid style={gray!25},
  legend style={font=\scriptsize,at={(0.5,1.03)},anchor=south,legend columns=2},table/col sep=comma]
\addplot[fill=omThm!60,draw=omThm] table[x=bin,y=squeeze]{figdata/strategies-1/08-short-interest-borrow-and-crowding/squeeze.csv};
\addplot[fill=omDef!60,draw=omDef] table[x=bin,y=actual]{figdata/strategies-1/08-short-interest-borrow-and-crowding/squeeze.csv};
\legend{si la moitié de l'encours vendeur se couvre en cinq jours,réel sur les cinq jours suivants}
\end{axis}
\end{tikzpicture}

\omcaption{Le livre vend à découvert les vingt noms les plus encombrés de la synthèse ($-0.5$ du capital)~: son rendement réel sur les cinq jours suivant chaque fin de mois, et sa perte
dans le scénario de squeeze, sur 96 fins de mois. Données~: \texttt{s1\_lending.squeeze}.}\label{fig:s1:short-interest-borrow-and-crowding:squeeze}
\end{omfigure}

\section{Fiches de stratégie}

\begin{strategyfile}{Évitement de l'encours vendeur élevé}\label{strat:s1:short-interest-borrow-and-crowding:avoid}
\sfield{Qui fournit le rendement, et pourquoi} Les vendeurs à découvert informés, dont les positions révèlent de l'information~; l'investisseur à l'achat seulement en profite sans rien payer.
\sfield{Instruments et plateformes} Actions, à l'achat seulement.
\sfield{Signal} L'encours vendeur rapporté au nombre d'actions ou au flottant, tel que publié en dernier lieu.
\sfield{Dimensionnement et exécution} Exclure les noms les plus vendus à découvert d'un portefeuille existant, ou les sous-pondérer.
\sfield{Coûts} Aucun coût propre, hormis celui du rééquilibrage.
\sfield{Comment elle meurt} Lorsque l'encours vendeur cesse d'être informatif~; les preuves pondérées par la valeur sont faibles.
\sfield{Horizon, capacité, infrastructure} Des mois~; une capacité très importante~; le fichier publié de l'encours vendeur.
\sfield{Backtest honnête} L'encours vendeur à sa date de publication, et non à sa date de règlement-livraison~; pondérations égales et pondérations par la valeur toutes deux.
\sfield{Sources} Asquith, Pathak et Ritter (2004/2005)~; la simulation de ce chapitre (0.94~\% par an, $t = 7.2$, un effet introduit à dessein).
\end{strategyfile}

\begin{strategyfile}{Signal du coût d'emprunt}\label{strat:s1:short-interest-borrow-and-crowding:fee}
\sfield{Qui fournit le rendement, et pourquoi} Une demande vendeuse supérieure à l'offre signale une surévaluation que les vendeurs ne peuvent pas entièrement corriger.
\sfield{Instruments et plateformes} Actions, avec un flux de prêt de titres.
\sfield{Signal} Les frais d'emprunt ou le taux d'utilisation, quotidiens.
\sfield{Dimensionnement et exécution} En signal d'évitement ou de sous-pondération~; vendre à découvert les noms eux-mêmes revient à payer les frais.
\sfield{Coûts} Les frais eux-mêmes~: sur le marché synthétique, ils dépassent la sous-performance.
\sfield{Comment elle meurt} Du côté vendeur, par construction~; rappels et pics de frais.
\sfield{Horizon, capacité, infrastructure} De quelques semaines à quelques mois~; données commerciales de prêt.
\sfield{Backtest honnête} Frais facturés chaque jour où une position est vendue à découvert, au taux alors en vigueur~; rappels simulés.
\sfield{Sources} Engelberg, Reed et Ringgenberg (2018)~; la simulation de ce chapitre.
\end{strategyfile}

\begin{strategyfile}{Classement par jours de couverture}\label{strat:s1:short-interest-borrow-and-crowding:dtc}
\sfield{Qui fournit le rendement, et pourquoi} Comme pour l'encours vendeur, pondéré par la difficulté à sortir de la position.
\sfield{Instruments et plateformes} Actions.
\sfield{Signal} L'encours vendeur divisé par le volume quotidien moyen.
\sfield{Dimensionnement et exécution} Acheter les titres à faible ratio, vendre à découvert ceux à ratio élevé, ou éviter.
\sfield{Coûts} Les frais d'emprunt du côté vendeur sont ici les plus élevés.
\sfield{Comment elle meurt} Par les squeezes sur les noms qu'elle vend à découvert.
\sfield{Horizon, capacité, infrastructure} De quelques semaines à quelques mois.
\sfield{Backtest honnête} Volume et encours vendeur aux mêmes dates~; frais et rappels.
\sfield{Sources} Boehmer, Jones et Zhang (2008)~: les actions fortement vendues à découvert ont sous-performé de 1.16~\% ajusté du risque sur 20 jours, en 2000--2004.
\end{strategyfile}

\begin{strategyfile}{Surcouche de risque des ventes encombrées}\label{strat:s1:short-interest-borrow-and-crowding:overlay}
\sfield{Qui fournit le rendement, et pourquoi} Personne~: c'est une assurance, financée en renonçant à une partie de la prime liée à l'encours vendeur.
\sfield{Instruments et plateformes} Les livres vendeurs de la société.
\sfield{Signal} Le score d'encombrement de chaque position vendeuse.
\sfield{Dimensionnement et exécution} Plafonner l'exposition du livre aux ventes encombrées~; limiter chaque nom par ses jours de couverture~; acheter des calls sur les noms les plus encombrés.
\sfield{Coûts} La prime abandonnée et les options achetées le cas échéant.
\sfield{Comment elle meurt} Elle ne meurt pas~; elle coûte.
\sfield{Horizon, capacité, infrastructure} Suivi quotidien~; données de prêt.
\sfield{Backtest honnête} Tester sous contrainte avec des scénarios de squeeze, et non seulement sur l'historique, qui contient rarement la queue de distribution.
\sfield{Sources} Rapport du personnel de la SEC (2021) sur GameStop~; le scénario de squeeze de ce chapitre.
\end{strategyfile}

\section{Tutoriel~: janvier 2021}\label{tut:s1:short-interest-borrow-and-crowding:tut}

\begin{tutorial}
\textbf{Objectif.} Ajouter un marché de prêt aux actions synthétiques, mesurer les signaux d'encours vendeur et de frais avant et après frais, et soumettre à un stress un livre vendeur encombré.
\textbf{État final~:} le tableau, la \cref{fig:s1:short-interest-borrow-and-crowding:book} et la \cref{fig:s1:short-interest-borrow-and-crowding:squeeze}.
\begin{tutsteps}
\item \textbf{Le marché de prêt}~: offre, demande informée et non informée, frais et rappels.
\omcode{code/firm/lendsig/firm_lendsig.py}{48}{75}{La couche de prêt.}\label{lst:s1:short-interest-borrow-and-crowding:lending}
\item \textbf{Le livre}~: un livre par déciles avec des frais sur ses positions vendeuses.
\omcode{code/strategies-1/08-short-interest-borrow-and-crowding/python/s1_lending.py}{86}{94}{Frais d'emprunt facturés du côté vendeur.}\label{lst:s1:short-interest-borrow-and-crowding:book}
\item \textbf{Exécuter} \texttt{levels}, \texttt{ic}, \texttt{book}, \texttt{avoidance}, \texttt{squeeze} et \texttt{fig\_lending.py}.
\end{tutsteps}
\textbf{Pour aller plus loin.} Faire en sorte que les rappels clôturent les positions et facturer leur réentrée~; laisser les vendeurs non informés se regrouper sur les noms informés (effet de troupeau)~; modéliser le squeeze
par un terme de rétroaction où la couverture fait monter le cours et la hausse du cours incite à couvrir.
\end{tutorial}

\section{Construction~: le marché de prêt}\label{bld:s1:short-interest-borrow-and-crowding:lendsig}

\begin{build}
\bfield{Rôle} Un marché de prêt de titres synthétique et les signaux et stress de vente à découvert qui en ont besoin.
\bfield{Interface} \texttt{LendingConfig}, \texttt{simulate\_lending(alpha, listed, shares, cfg)}, \texttt{fee\_of(util, cfg)},
\texttt{days\_to\_cover(si, shares, adv\_shares)}, \texttt{crowding(si, util, dtc)}, \texttt{squeeze\_loss(w, si, shares, adv\_shares, sigma, cover, days, eta)}.
\bfield{Règles} L'encours vendeur ne dépasse jamais l'offre~; frais facturés chaque jour sur les ventes à découvert~; rappels aléatoires, avec une probabilité croissant avec le taux d'utilisation.
\bfield{Tests d'acceptation} \texttt{code/firm/lendsig/tests/}~: le barème des frais calculé à la main~; la demande informée concentrée sur les noms à alpha négatif~; les jours de couverture,
le score d'encombrement et la perte de squeeze calculés à la main.
\bfield{Pour aller plus loin} La réhypothécation (encours vendeur supérieur à 100~\% du flottant)~; l'offre de prêt venant des fonds indiciels~; un squeeze avec rétroaction.
\end{build}

\begin{omsources}
\item US SEC staff, \emph{Staff Report on Equity and Options Market Structure Conditions in Early 2021}, octobre 2021.
\item P.~Asquith, P.~A.~Pathak et J.~R.~Ritter, ``Short interest, institutional ownership, and stock returns'', \emph{Journal of Financial Economics}
78(2), 2005.
\item E.~Boehmer, C.~M.~Jones et X.~Zhang, ``Which shorts are informed?'', \emph{Journal of Finance} 63(2), 2008.
\item J.~E.~Engelberg, A.~V.~Reed et M.~C.~Ringgenberg, ``Short-selling risk'', \emph{Journal of Finance} 73(2), 2018.
\item FINRA Rule 4560, déclaration de l'encours vendeur à découvert.
\end{omsources}

\section{Exercices}

\begin{exercise}[$\star$]\label{exo:s1:short-interest-borrow-and-crowding:1}
Un titre compte 80 millions d'actions en circulation, dont 10~\% sont vendues à découvert, et se traite à 0.8 million d'actions par jour. Calculer ses jours de couverture.
\end{exercise}

\begin{exercise}[$\star$]\label{exo:s1:short-interest-borrow-and-crowding:2}
Comment l'encours vendeur peut-il dépasser 100~\% du flottant~? Donner un exemple en deux étapes.
\end{exercise}

\begin{exercise}[$\star$]\label{exo:s1:short-interest-borrow-and-crowding:3}
Avec le barème de frais du chapitre, combien coûte l'emprunt d'une action à 80~\% d'utilisation, et combien coûte une position vendeuse de 5~\% du capital sur ce titre
pendant neuf mois~?
\end{exercise}

\begin{exercise}[$\star\star$]\label{exo:s1:short-interest-borrow-and-crowding:4}
Pourquoi le livre des frais perd-il alors que le livre de l'encours vendeur gagne, lorsque les deux visent les mêmes noms~?
\end{exercise}

\begin{exercise}[$\star\star$]\label{exo:s1:short-interest-borrow-and-crowding:5}
Pourquoi la version d'évitement du signal est-elle plus robuste que la version longue--courte~?
\end{exercise}

\begin{exercise}[$\star\star$]\label{exo:s1:short-interest-borrow-and-crowding:6}
Un livre est vendeur de vingt noms encombrés à parts égales, avec la moitié de son capital. L'un d'eux est multiplié par vingt. Que perd le livre~?
\end{exercise}

\begin{exercise}[$\star\star\star$]\label{exo:s1:short-interest-borrow-and-crowding:7}
\emph{Programmation.} Exécuter \texttt{squeeze(0.25, 5)} et \texttt{squeeze(0.5, 10)}. Expliquer comment les pertes dépendent de la part couverte et de la vitesse de
couverture sous la loi en racine carrée.
\end{exercise}

\begin{exercise}[$\star\star\star$]\label{exo:s1:short-interest-borrow-and-crowding:8}
\emph{Trouver la faille.} «~Notre backtest de l'encours vendeur utilise l'encours publié par la FINRA à sa date de règlement-livraison et rapporte 1.5~\% par mois.~»
\end{exercise}

\section{Problème~: janvier 2021}

\begin{problem}[{Problème du week-end --- le signal et le squeeze}]\label{pb:s1:short-interest-borrow-and-crowding:1}
Le marché synthétique du prêt et le registre public.

\medskip\noindent\textbf{Partie I --- L'épisode.}
\begin{enumerate}
\item À quel niveau l'encours vendeur de GameStop a-t-il atteint, et comment peut-il dépasser 100~\%~?
\item Décrire l'évolution du cours du 11 au 28 janvier 2021.
\item À quelle conclusion le personnel de la SEC est-il parvenu au sujet de la couverture des positions vendeuses~?
\item Comment l'encours vendeur est-il publié, et à quel point est-il daté~?
\end{enumerate}

\medskip\noindent\textbf{Partie II --- Le signal.}
\begin{enumerate}[resume]
\item Que mettent en évidence Boehmer, Jones et Zhang, ainsi qu'Asquith, Pathak et Ritter~?
\item Décrire le marché synthétique du prêt.
\item Donner les ICs de l'encours vendeur, des jours de couverture, des frais et du score d'encombrement.
\item Donner les rendements et les ratios de Sharpe du livre de l'encours vendeur, avant et après frais.
\end{enumerate}

\medskip\noindent\textbf{Partie III --- Les coûts.}
\begin{enumerate}[resume]
\item Pourquoi le livre des frais perd-il~?
\item Que rapporte l'évitement~?
\item Qu'Engelberg, Reed et Ringgenberg ont-ils mis en évidence au sujet du risque de la vente à découvert~?
\item Comment les rappels sont-ils modélisés, et pourquoi se concentrent-ils sur les noms du signal~?
\end{enumerate}

\medskip\noindent\textbf{Partie IV --- Le verdict.}
\begin{enumerate}[resume]
\item Énoncer le \emph{résultat nommé}~: le rendement du signal d'encours vendeur élevé avant et après frais d'emprunt, et la perte d'un livre vendeur encombré dans un
squeeze simulé.
\item Pourquoi les pertes de squeeze sont-elles des planchers et non des prévisions~?
\item Que fait un nom multiplié par vingt à un livre encombré~?
\item Comment faudrait-il dimensionner les noms encombrés d'un livre vendeur~?
\item Quelle fiche de stratégie un gérant en position longue uniquement utiliserait-il~?
\item Quelles données faudrait-il acheter pour faire tourner ces stratégies~?
\item Qu'enseigne l'épisode GameStop sur le risque de queue dans les backtests~?
\item En une phrase~: que paie un vendeur à découvert~?
\end{enumerate}
\end{problem}

\section{Questions d'entretien}

\begin{interviewq}[$\star$ \iqroles{chercheur, trader}]\label{iq:s1:short-interest-borrow-and-crowding:1}
Que vous dit un encours vendeur élevé sur une action~?
\end{interviewq}

\begin{interviewq}[$\star\star$ \iqroles{trader, risque}]\label{iq:s1:short-interest-borrow-and-crowding:2}
Quels risques une position vendeuse porte-t-elle et qu'une position longue ne porte pas~?
\end{interviewq}

\begin{interviewq}[$\star\star$ \iqroles{risque}]\label{iq:s1:short-interest-borrow-and-crowding:3}
Comment mesureriez-vous et limiteriez-vous l'encombrement d'un livre vendeur~?
\end{interviewq}

\begin{interviewq}[$\star\star$ \iqroles{chercheur}]\label{iq:s1:short-interest-borrow-and-crowding:4}
Pourquoi un signal qui prédit une sous-performance pourrait-il encore perdre de l'argent lorsqu'on le négocie du côté vendeur~?
\end{interviewq}

\begin{interviewq}[$\star\star$ \iqroles{développeur, chercheur}]\label{iq:s1:short-interest-borrow-and-crowding:5}
Quels horodatages un backtest de l'encours vendeur exige-t-il, et qu'est-ce qui tourne mal sans eux~?
\end{interviewq}

\begin{interviewq}[$\star\star\star$ \iqroles{chercheur, risque}]\label{iq:s1:short-interest-borrow-and-crowding:6}
Les vendeurs à découvert détiennent $Q$ actions d'un titre qui s'échange à $V$ par jour, avec une volatilité quotidienne $\sigma$. Sous l'impact en racine carrée, comment le cours évolue-t-il s'ils couvrent tous en
$d$ jours, et pourquoi le modèle sous-estime-t-il un squeeze~?
\end{interviewq}
